% !PRINTED 631
\weberpage{631}
Machen wir darin die Substitution $(r, r^n)$, so wird
\weq{14}{r^n \equiv g^{-\frac{p-1}{m}} \ (\mathrm{mod}\ \mathfrak p_n),}
oder wenn wir $n'$ durch die Congruenz
\weq{15}{n n' \equiv 1 \ (\mathrm{mod}\ m)}
definiren:
\weq{16}{r \equiv g^{-n'\frac{p-1}{m}} \ (\mathrm{mod}\ \mathfrak p_n).}

Es ist also in diesem Falle, übereinstimmend mit der allgemeinen
 Regel, $\mathfrak p_n$ der grösste gemeinschaftliche Theiler von
\[
p \ \text{und} \ \left(r - g^{-n'\frac{p-1}{m}}\right).
\]

Durch diese Bestimmung sind die einzelnen Primformen $\mathfrak p_n$
genau charakterisirt. Wir erhalten also den Satz:

2. \gesperrt{Eine Primzahl $p$, die nach dem Modul $m$ mit 1
congruent ist, zerfällt im Körper $\Omega_m$ in $\varphi(m)$ von
einander verschiedene Primfactoren $\mathfrak p_n$ vom
ersten Grade, die man als grösste gemeinschaftliche
 Theiler von $p$ mit den verschiedenen
Zahlen $r - g^{-n'\frac{p-1}{m}}$ erhält, wenn $g$ eine primitive
Wurzel von $p$ ist und $n'$ durch die Congruenz
$n n' \equiv 1 \ (\mathrm{mod}\ m)$ bestimmt ist.}

Nach \S.\,141 können wir jede ganze oder gebrochene Zahl $\omega$
und jedes Functional des Körpers $\Omega_m$ in der Weise in Primfactoren
 zerlegen, dass Zähler und Nenner keinen gemeinschaftlichen
 Primfactor enthalten, und diese Zerlegung ist, von
Einheitsfactoren abgesehen, völlig bestimmt, so dass, wenn
$\mathfrak p$ ein Primideal ist, ein ganz bestimmter positiver oder negativer
Exponent $k$ existirt, so dass das Functional $\omega \mathfrak p^{-k}$ den Primfactor
 $k$ weder im Zähler noch im Nenner enthält. % [sic] printed $k$; context requires $\mathfrak p$
Wir sagen
dann der Kürze halber, es sei $\mathfrak p^k$ \gesperrt{die in $\omega$ enthaltene Potenz
von $\mathfrak p$.}

Ist $k=0$, so sagen wir, $\mathfrak p$ ist in $\omega$ \gesperrt{nicht enthalten}.

Zwei Zahlen (oder auch Functionale) $\omega, \omega'$ eines Körpers, in
denen ein und dasselbe Primfunctional $\mathfrak p$ enthalten ist, \gesperrt{wollen
wir theilerverwandt oder kurz verwandt nennen}; wenn dagegen
kein Primfunctional zugleich in $\omega$ und in $\omega'$ enthalten ist, so
heissen $\omega$ und $\omega'$ \gesperrt{theilerfremd oder fremd}. Besonders nützlich
wird uns dieser Ausdruck in der Folge sein, wenn an
%CONT>>
